\documentclass[11pt]{scrartcl}

\usepackage[stixtwo]{fontsetup}
\usepackage{amsmath}

\begin{document}
Teljes indukcióval megmutatjuk, hogy $1 \le d_n \le 6$ minden
$n \in \mathbb{N}^+$ esetén. Mivel $d_1 = 2$, ezért $n = 1$-re az állítás
teljesül. Ha valamilyen $n$-re teljesül, azaz $1 \le d_n \le 6$, akkor
$\frac{1}{6} \le \frac{1}{d_n} \le 1$, így $1 \le \frac{6}{d_n} \le 6$, és
ezért $1 \le 7 - \frac{6}{d_n} = d_{n + 1} \le 6$.

Teljes indukcióval megmutatjuk, hogy $d_n$ monoton növő, azaz
$d_n \le d_{n + 1}$ minden $n \in \mathbb{N}^+$ esetén. Mivel $d_2 = 4$,
ezért $n = 1$-re az állítás teljesül. Ha valamilyen $n$-re teljesül, azaz
$d_n \le d_{n + 1}$, akkor mivel $d_n$ pozitív, ezért
$\frac{1}{d_n} \ge \frac{1}{d_{n + 1}}$, és így
\[
	d_{n + 1} = 7 - \frac{6}{d_n} \le 7 - \frac{6}{d_{n + 1}} = d_{n + 2},
\]
vagyis $n + 1$-re is teljesül az egyenlőtlenség.

Mivel a sorozat korlátos és monoton, ezért konvergens, határértéke a
rekurzió valamelyik fixpontja. Az $A = 7 - \frac{6}{A}$ egyenlet megoldásai
$A = 1$ és $A = 6$. Mivel $d_n$ monoton növő és $d_2 = 4$, a határérték nem
lehet $1$. Tehát $\lim d_n = 6$.
\end{document}
